% Fichier généré automatiquement par tools/generate_corrections.py.
% Source : CIN/CIN-03-Transmetteurs/25_Cheville/25_Cheville.tex
% Statut : complété (6/6 question(s)).
\begin{corrigebox}[Corrigé complété]
\CorrigeQuestion{1}
D'après le diagramme de définition des blocs et le diagramme des exigences, les rapports de transmission doivent être : 
\begin{itemize}
\item pour l'axe de tangage : $\dfrac{N_{\text{moteur}}}{N_{\text{Tangage}}}=138,33$ au minimum; 
\item pour l'axe de roulis :  $\dfrac{N_{\text{moteur}}}{N_{\text{Roulis}}}= 197,61$ au minimum.
\end{itemize}
\CorrigeQuestion{2}
\begin{center}
\begin{tabular}{lccc}
\hline
\textbf{Roue dentée} & \textbf{Module} & \textbf{Nb dents} & \textbf{Diamètre (mm)}\\
\hline
Pignon 03 20 & 0,3 &20        & 6		\\ 
Mobile Inf1 Roue & 0,3 & 80  & 24 	\\ 
Mobile Inf1 Pignon & 0,4 & 25 & 10	\\ 
Mobile Inf2 Roue & 0,4 & 47   & 18,8	\\ 
Mobile Inf2 Pignon & 0,4 & 12 & 4,8	\\ 
Mobile Inf4 Roue & 0,4 & 58   & 23,2	\\ 
Mobile Inf4 Pignon & 0,7 & 10 & 7	\\  
Roue de sortie & 0,7 & 36      & 25,2 \\ \hline
\end{tabular}
\end{center}
\CorrigeQuestion{3}
On choisit les inconnues et les conventions positives de la figure,
        on écrit les relations de conservation/fermeture adaptées, puis on résout le
        système obtenu. Le résultat symbolique doit être homogène et vérifié dans les cas
        limites (entrée nulle, rapport unitaire ou position de référence).
\CorrigeQuestion{4}
Le schéma est reconstruit en conservant uniquement les axes, centres,
        directions de glissière et longueurs fonctionnelles. Les angles et courses imposés
        sont appliqués dans l'ordre de la chaîne cinématique; les liaisons doivent rester
        fermées et les contacts sans pénétration.
\CorrigeQuestion{5}
$$
R_T = (-1)^n \dfrac{80\cdot 47 \cdot 58 \cdot 36}{20\cdot 25\cdot 12 \cdot 10 } = 130,85
$$

Ceci est inférieur à ce qui est préconisé par le cahier des charges. 

Pour respecter le cahier des charges, on peut :
\begin{itemize}
\item choisir un autre moteur;
\item changer le nombre de dents d'une des roues. Il suffirait pour cela que,  par exemple, la roue de sortie comporte 39 dents. 
\end{itemize}
\CorrigeQuestion{6}
Le rapport de transmission du second train est de 201,3 ce qui est compatible avec le cahier des charges.
\end{corrigebox}
